\section{Conclusion}\label{chap:conclusion}
This project explored Support Vector Machines (SVMs) theoretically and practically. We formulated the primal optimization problem, derived the dual formulation, and demonstrated how the kernel trick enables effective non-linear classification.

Key theoretical properties of SVMs include:
\begin{itemize}
    \item The convex nature of the primal problem guarantees a global optimum
    \item The dual formulation offers computational advantages, particularly with high-dimensional data
    \item Kernel functions enable implicit mapping to high-dimensional spaces without computational burden
    \item The sparse solution relies only on support vectors, increasing prediction efficiency
\end{itemize}

Our experiments confirmed these foundations, with primal and dual approaches converging to nearly identical solutions for linearly separable data. Regularization parameter experiments demonstrated the expected bias-variance trade-off. For non-linear data, RBF, polynomial, and Laplacian kernels achieved perfect classification where linear kernels failed, confirming the kernel trick's effectiveness.

Application to the Breast Cancer Wisconsin dataset yielded 98.25\% accuracy with an RBF kernel, demonstrating SVMs' value in medical diagnostics. PCA visualization showed tumor concavity features as important discriminators.

Despite their effectiveness, SVMs have limitations regarding optimization requirements, scaling issues with large datasets, and binary classification constraints.

Future work could include:
\begin{itemize}
    \item More efficient optimization techniques like Sequential Minimal Optimization
    \item Extension to multi-class classification
    \item Automatic kernel selection and parameter tuning methods
    \item Hybrid models combining SVMs with deep learning
\end{itemize}

In conclusion, this project demonstrated SVMs' mathematical elegance and practical utility. Their theoretical guarantees and strong empirical performance justify their continued relevance despite newer machine learning approaches.
